\subsection{The norm count}

For unit-amplitude nearest-neighbour hopping across one bond,
\[
 K_{xy}=\sum_{\alpha,\beta=1}^{16}
 \left(a_{x,\alpha}^\dagger U_{\alpha\beta}a_{y,\beta}
 +a_{y,\beta}^\dagger U_{\alpha\beta}^*a_{x,\alpha}\right),
\tag{6.4}\label{eq:hopping}
\]
with $U$ unitary, the one-particle bond matrix has sixteen eigenvalues
$+1$ and sixteen eigenvalues $-1$.  Its fermionic second quantization
therefore has
\[
 \norm{K_{xy}}=16.
\tag{6.5}\label{eq:bondnorm}
\]
If the lattice coordination is $z$, at most $z$ such bonds touch one
cluster.  The advertised count is consequently
\[
 c_{\mathrm{count}}=16z,\qquad
 t_{0,\mathrm{count}}=\frac1{32z}.
\tag{6.6}\label{eq:count}
\]
On a three-dimensional cubic spatial slice, $z=6$, so
\[
 c_{\mathrm{count}}=96,\qquad
 t_{0,\mathrm{count}}=\frac1{192}
 =0.005208333\ldots .
\tag{6.7}\label{eq:count3d}
\]

This count controls the norm incident on one site.  It does not prove
\eqref{eq:block} or \eqref{eq:formbound}.  In fact ordinary hopping
creates a pair of defective clusters from components of the product
ground state, so
\[
 QK P\ne0.
\tag{6.8}\label{eq:notblock}
\]
The numeric twin measures
$\norm{QKP}=\sqrt2$ on the two-cluster four-mode analogue.  Thus
Proposition~\ref{prop:relative} cannot be applied to
\eqref{eq:hopping}, and \eqref{eq:count3d} is not the explicit
$t_*$ of Theorem~\ref{thm:T3}.  A global Weyl estimate would only give
$1-2|t|\norm V$, whose right side scales with volume and is unsuitable.

\begin{center}
\setlength{\fboxsep}{9pt}
\fcolorbox{orange!75!black}{orange!4}{
\begin{minipage}{0.91\textwidth}
\textbf{Uncertified explicit-threshold box.}
The number $t_{0,\mathrm{count}}=1/192$ is the value obtained from the
requested local norm count on a cubic spatial slice.  It is a valid
threshold for the \emph{conditional} relative-form proposition if
\eqref{eq:block}--\eqref{eq:formbound} are separately proved with
$c=96$.  Ordinary hopping fails \eqref{eq:block}, visibly and
numerically.  The applicable Michalakis--Zwolak theorem instead gives
$t_*=J_0/J_V>0$ with $J_0=1/(3c_{\rm MZ})$, where $c_{\rm MZ}$ contains
the quasi-local spectral-flow decay $w$ in \eqref{eq:MZc}.  No value of
$c_{\rm MZ}$ has been derived for this model, so replacing it by $96$
would be an unjustified theorem claim.
\end{minipage}}
\end{center}

\section{Numeric twin}

The companion probe is
\[
 \texttt{experiments/tfpt-discovery/det16\_gap\_stability\_probe.py}.
\]
It uses the exact four-mode determinant projectors and open-chain
hopping machinery of the frozen \texttt{v1002} route.  All $19$ checks
pass.  The direct comparison is:
\[
\begin{array}{c|rrrr}
\toprule
\text{clusters}&t=0.05&t=0.10&t=0.20&t=0.30\\
\midrule
2&0.973726995&0.946980790&0.902548402&0.837138848\\
3&0.962836264&0.924929608&0.860366685&0.764340648\\
4&0.957577627&0.914546582&0.842484121&0.726091276\\
\bottomrule
\end{array}
\tag{7.1}\label{eq:numeric}
\]
Every entry lies above $1-2c_{\rm count}|t|$ with
$c_{\rm count}=4$ for the two-cluster endpoint pair and $8$ for the
three--four cluster chains.  Except at the smallest $t$, this linear
line is already negative and hence not tight.

The scan uses step $0.05$ for two and three clusters and step $0.10$
for four clusters on $0\leq t\leq4$.  It finds no finite-volume gap
closing.  The minimum scanned gaps are
\[
\begin{array}{c|cc}
\toprule
\text{clusters}&\min\gap&\text{minimizing }t\\
\midrule
2&0.764490286&0.55\\
3&0.472500856&4.00\\
4&0.549772255&0.60\\
\bottomrule
\end{array}
\tag{7.2}\label{eq:scan}
\]
Consequently the requested finite-chain ``actual closing''
$t_c$ is not located because it does not occur in the declared scan:
\[
 t_c^{\rm scan}>4,\qquad
 \frac{t_{0,\mathrm{count}}^{\rm toy}}{t_c^{\rm scan}}
 <\frac{1/16}{4}=0.015625.
\tag{7.3}\label{eq:margin}
\]
This is an upper bound on a ratio, not an invented critical point.

The positive and negative controls expose the theorem hypotheses:
\[
\begin{array}{ll}
\text{Spin(4)-covariant flavor-diagonal hopping:}
 &\max_A\norm{[K,G^A]}=0,\\
\text{flavor-selective mutant:}
 &\max_A\norm{[K_{\rm mut},G^A]}=1,\\
\text{$1/3$ normalization mutant:}
 &\norm{\widetilde P^2-\widetilde P}=1/9,\quad
 \Tr\widetilde P=2/3.
\end{array}
\tag{7.4}\label{eq:mutants}
\]
Thus one mutant visibly leaves the covariant perturbation class and the
other destroys the projector/frustration-free premise.

\section{Combined strongest theorem}

\begin{theorem}[Volume-independent DET16 mirror-gap theorem]
\label{thm:main}
For the matter-only DET16 cluster model
\eqref{eq:cluster}--\eqref{eq:H0} on any finite bounded-degree lattice:
\begin{enumerate}[label=\textup{(\alph*)}]
\item the projectors are mutually commuting even CAR operators;
$H_0$ has exact spectrum $\{0,\ldots,|\Lambda|\}$, multiplicities
\eqref{eq:multiplicity}, unique product ground state
\eqref{eq:productground}, and gap one in every volume;
\item $H_0$ is locally Spin(10)-Gauss invariant, parity invariant, and
translation invariant for constant $\phi$; its unique product ground
state has neither spontaneous symmetry breaking nor torus/topological
degeneracy;
\item every uniformly bounded finite-range even perturbation has a
volume-independent stability interval $|t|\leq t_*$ on which the gap
is at least $1/2$, with $t_*>0$ supplied by the
Michalakis--Zwolak theorem after the exact checks
$\Delta_0=0$ and $\gamma(r)=1$;
\item for finite $\mathbb Z_2/\mathbb Z_4$ links on an open chain,
$h_E>0$, and gauge-covariant nearest-neighbour hopping, the physical
Gauss sector has a unique ground state and volume-independent gap at
least $\Delta/2$ throughout the cited window
$|t|<\Delta\tau_{\rm FP}/16$, with
$\Delta=\min(1,h_E)$ and $\tau_{\rm FP}>0$ parameterized by
\eqref{eq:FPparametric};
\item if, additionally, a perturbation obeys the block-diagonal
relative-form hypotheses \eqref{eq:block}--\eqref{eq:formbound}, then
the explicit elementary estimate \eqref{eq:linearbound} holds.
\end{enumerate}
No numerical identification of $t_*$ with
$t_{0,\mathrm{count}}=1/(32z)$ is asserted.
\end{theorem}

\begin{proof}
Part (a) is Theorem~\ref{thm:T1} and Lemma~\ref{lem:order}.
Part (b) is Theorem~\ref{thm:T2}.  Part (c) is
Theorem~\ref{thm:T3}, whose proof verifies every hypothesis of the
quoted stability theorem against the concrete product model.  Part
(d) is Proposition~\ref{prop:FPaudit}, including the ambient-to-physical
gap restriction.  Part (e) is Proposition~\ref{prop:relative}.  The
final sentence follows from \eqref{eq:notblock} and the
explicit-threshold box.
\end{proof}

\section{Physical spectator and remaining boundary}
\label{sec:boundary}

\subsection{Exact spectator statement}

Let $\Hh_{\rm wall}$ carry a wall Hamiltonian $H_{\rm wall}$ and let
$\Hh_{\rm mir}$ carry $H_{0,\Lambda}$.  If the DET16 terms act as the
identity on wall modes, then
\[
 H_{\rm total}(0)
 =H_{\rm wall}\otimes\id_{\rm mir}
 +\id_{\rm wall}\otimes H_{0,\Lambda}.
\tag{9.1}\label{eq:spectator}
\]
Hence
\[
 \operatorname{spec}H_{\rm total}(0)
 =\{E_{\rm wall}+k:
 E_{\rm wall}\in\operatorname{spec}H_{\rm wall},\
 k=0,\ldots,|\Lambda|\}.
\tag{9.2}\label{eq:spectralsum}
\]
The $k=0$ copy is exactly the original wall spectrum, with every
multiplicity preserved.

For $t\ne0$, ``does not mix'' must mean that no interaction term
contains both wall and mirror operators.  Exact preservation, rather
than merely perturbative continuity, further requires
\[
 V=\id_{\rm wall}\otimes V_{\rm mir}.
\tag{9.3}\label{eq:nomix}
\]
Then the total Hamiltonian remains a tensor sum, and the stable mirror
ground band adds one common scalar energy to the wall spectrum.  If a
wall-only term is allowed it changes the wall spectrum on its own; if a
wall--mirror coupling is allowed, \eqref{eq:spectralsum} is no longer
exact and a separate relative/local stability estimate is required.

\begin{center}
\setlength{\fboxsep}{10pt}
\fcolorbox{red!65!black}{red!3}{
\begin{minipage}{0.91\textwidth}
\textbf{Remaining boundary for
\texttt{CHIRAL4D.MIRROR.DET16.01}.}

\smallskip
\textbf{Compact/continuous dynamical gauge fields.}
Section~\ref{sec:dynlinks} closes the finite
$\mathbb Z_2/\mathbb Z_4$ open-chain item by
Theorem~\ref{thm:FP} and Proposition~\ref{prop:FPaudit}; it is no longer
part of this boundary box.  Compact-group rotor links have
infinite-dimensional local Hilbert spaces and therefore violate the
finite-$M$ hypothesis of the bounded-interaction FP theorem quoted here.
They require the form-bounded extension with its domain hypotheses, or
truncation-uniform estimates, neither of which is audited in this
memorandum.

\smallskip
\textbf{Domain wall.}
At $t=0$ the spectator theorem
\eqref{eq:spectator}--\eqref{eq:spectralsum} is exact.  For
$|t|<t_*$, exact preservation of the wall spectrum requires the
factorization \eqref{eq:nomix}.  Gauge covariance or fermion parity
alone does not forbid wall--mirror bilinears, so the microscopic
domain-wall construction must prove their absence or bound them.

\smallskip
\textbf{Three spatial dimensions.}
The count $z=6$ and $1/192$ assumes a cubic three-dimensional spatial
slice, but no lattice domain-wall Hamiltonian, chirality localization,
continuum limit, or Lorentz-covariant $3+1$-dimensional theory is
constructed here.  The theorem is a uniform quantum-lattice statement
for the mirror matter sector, not yet the physical four-dimensional
placement.

\smallskip
\textbf{Explicit hopping threshold.}
The applicable stability theorem proves $t_*>0$ and gap at least
$1/2$, but its constant contains the spectral-flow decay
$c_{\rm MZ}$.  The simple value $1/192$ is conditional and cannot be
promoted to the hopping theorem because $QKP\ne0$.  Deriving a
numerical $c_{\rm MZ}$, or proving a different explicit local
block-diagonalization/Knabe estimate, remains open.
\end{minipage}}
\end{center}

\section*{Conclusion}
\addcontentsline{toc}{section}{Conclusion}

T1 and T2 are elementary, exact, and volume independent.  T3 is a
genuine thermodynamic stability statement: the product projector has
zero Local-TQO error and unit local gap, so the full
Michalakis--Zwolak theorem applies to bounded finite-range even
perturbations.  Its threshold is uniform but not numerically evaluated.
The requested count $1/(32z)$ is retained only in the precise
conditional proposition it actually satisfies; the numeric twin both
passes the comparison and detects the missing block-diagonal premise.
The dynamical-link audit repairs the tensor-product endpoint exactly and
applies the Fröhlich--Pizzo theorem, which does control the spectral gap
uniformly in chain length.  This moves finite
$\mathbb Z_2/\mathbb Z_4$ open chains to cited-verified status with the
parametric window $|t|<\Delta\tau_{\rm FP}/m$; it does not manufacture a
numerical value for the paper's implicit induction constant.  Accordingly,
the red box now retains only compact/continuous links, domain-wall
mixing, and the $3+1$-dimensional placement, preventing the chain theorem
from being mistaken for the complete dynamical-gauge contract.

\begin{thebibliography}{9}

\bibitem{BHM2010}
S.~Bravyi, M.~B.~Hastings, and S.~Michalakis,
\emph{Topological quantum order: stability under local perturbations},
J.\ Math.\ Phys.\ \textbf{51} (2010), 093512,
\href{https://doi.org/10.1063/1.3490195}
{doi:10.1063/1.3490195}.

\bibitem{MZ2013}
S.~Michalakis and J.~P.~Zwolak,
\emph{Stability of Frustration-Free Hamiltonians},
Commun.\ Math.\ Phys.\ \textbf{322} (2013), 277--302,
\href{https://doi.org/10.1007/s00220-013-1762-6}
{doi:10.1007/s00220-013-1762-6}.

\bibitem{BDL2011}
S.~Bravyi, D.~P.~DiVincenzo, and D.~Loss,
\emph{Schrieffer--Wolff transformation for quantum many-body systems},
Ann.\ Phys.\ \textbf{326} (2011), 2793--2826,
\href{https://doi.org/10.1016/j.aop.2011.06.004}
{doi:10.1016/j.aop.2011.06.004}.

\bibitem{FP2020}
J.~Fröhlich and A.~Pizzo,
\emph{Lie--Schwinger Block-Diagonalization and Gapped Quantum Chains},
Commun.\ Math.\ Phys.\ \textbf{375} (2020), 2039--2069,
\href{https://doi.org/10.1007/s00220-019-03613-2}
{doi:10.1007/s00220-019-03613-2};
\href{https://arxiv.org/abs/1812.02457}
{arXiv:1812.02457}.

\end{thebibliography}

\end{document}
